\chapter{A Concrete Naming Certificate for $\SubjectReductionRouteEquivalenceUp$}
\label{ch:concrete-instances-subject-reduction-route-equivalence-namecert}
\origin{human}

The $\SubjectReductionRouteEquivalenceUp$ packet gives the route vision of \autoref{ch:sr-discharge-routes} a finite BEDC landing. It reads the socket rows of \autoref{ch:concrete-instances-subject-reduction-discharge-socket-namecert} and sits beside the classifier packet of \autoref{ch:concrete-instances-subject-reduction-route-classifier-namecert}: the classifier names the explicit bundle route being used, while the equivalence packet records that every accepted consumer-facing discharge read terminates at that bundle surface.

\begin{definition}[Subject-reduction route equivalence carrier]
\label{def:subject-reduction-route-equivalence-carrier}
A $\SubjectReductionRouteEquivalenceUp$ carrier is a finite $\BHist$ packet
$$
S=(B,D,H,C,P,N)
$$
with the following displayed rows:
\begin{enumerate}
\item $B$ is the bundle-route row, recording explicit construction of the four-field subject-reduction discharge bundle;
\item $D$ is the bundle-consumption row, recording the finite handoff from the bundle row $B$ to the socket surface;
\item $H$ gives componentwise $\hsame$ transport for the four discharge coordinates, the bundle row, the consumer route, provenance, and local naming;
\item $C$ gives the displayed $\Cont$ consumer routes through which downstream packets read the route agreement;
\item $P$ gives $\Pkg$ provenance over $\BHist$, $\Ext$, $\Cont$, $\SubjectReductionDischargeSocketUp$, $\SubjectReductionRouteClassifierUp$, the bundle route, componentwise $\hsame$, and $\NameCert$ rows;
\item $N$ is the local $\NameCert_{\SubjectReductionRouteEquivalenceUp}$ row.
\end{enumerate}
The carrier contains no inhabitant of the four discharge hypotheses, no closed subject-reduction theorem, and no route equality outside the displayed finite rows.
\end{definition}

\begin{definition}[Subject-reduction route equivalence classifier]
\label{def:subject-reduction-route-equivalence-classifier}
The $\SubjectReductionRouteEquivalenceUp$ classifier compares only the rows displayed in \autoref{def:subject-reduction-route-equivalence-carrier}. Two accepted packets agree exactly when their bundle rows $B$, consumption rows $D$, componentwise $\hsame$ ledgers $H$, consumer-route rows $C$, provenance rows $P$, and local naming rows $N$ agree through the displayed $\BHist$ coordinates.

The classifier has no coordinate for host definitional equality between route implementations, no proof that the route is inhabited for the closed-CIC substrate, and no hidden discharge witness. Its entire comparison surface is the finite packet $S=(B,D,H,C,P,N)$.
\end{definition}

\begin{theorem}[SubjectReductionRouteEquivalence bundle-socket handoff]
\label{thm:subject-reduction-route-equivalence-bundle-socket-handoff}
For every accepted $\SubjectReductionRouteEquivalenceUp$ carrier $S=(B,D,H,C,P,N)$, the bundle route $B$ feeds the $\SubjectReductionDischargeSocketUp$ surface through $D$. A consumer reaching the socket through $B$ reads the four displayed discharge coordinates in bundle order through the componentwise transport rows in $H$.
\end{theorem}
\begin{proof}
Project the carrier rows. The row $B$ already has the bundle-shaped socket order.

Read the consumption row $D$. It is the only permitted route from $B$ to the socket-facing bundle shape, so any consumer must enter the socket through that displayed handoff.

Finally use the transport rows $H$ and the consumer routes $C$. They identify the four coordinates componentwise and replay them to the socket surface. Since the packet has no further route coordinate, every read terminates at the same displayed socket rows.
\end{proof}

\begin{theorem}[SubjectReductionRouteEquivalence NameCert obligations]
\label{thm:subject-reduction-route-equivalence-namecert-obligations}
For the carrier of \autoref{def:subject-reduction-route-equivalence-carrier}, the local $\NameCert_{\SubjectReductionRouteEquivalenceUp}$ surface consists exactly of five obligations:
\begin{enumerate}
\item carrier admission for the finite packet $S=(B,D,H,C,P,N)$;
\item row classification for the bundle row $B$ and consumption row $D$;
\item componentwise stability through the $\hsame$ transport ledger $H$;
\item consumer exactness through the displayed $\Cont$ routes $C$;
\item package provenance and local naming through $P$ and $N$.
\end{enumerate}
\end{theorem}
\begin{proof}
Carrier admission is the first projection of the displayed packet $S$. The next projection reads the route rows $B,D$ and therefore supplies the complete classifier surface for bundle coordinates.

The transport obligation is exactly the row $H$. It compares the four discharge coordinates and the route metadata componentwise, without adding another discharge source.

The remaining public obligations are the consumer routes and provenance rows. The row $C$ gives the only consumer reads, while $P$ and $N$ bind package origin and local naming. These five groups exhaust the carrier coordinates.
\end{proof}

\begin{theorem}[SubjectReductionRouteEquivalence consumer non-escape]
\label{thm:subject-reduction-route-equivalence-consumer-nonescape}
Every consumer read of an accepted $\SubjectReductionRouteEquivalenceUp$ carrier receives route agreement only as the finite rows
$$
B,T,D,H,C,P,N.
$$
The consumer may use these rows to compare bundle-route reads at the $\SubjectReductionDischargeSocketUp$ surface. It cannot infer an inhabited discharge bundle, a closed subject-reduction theorem, host equality of route implementations, or any discharge evidence not present in the packet.
\end{theorem}
\begin{proof}
Apply the NameCert-obligation theorem to identify every accepted read with one of the displayed carrier rows.

For a permitted route-agreement read, the consumer follows $B$ directly to the socket or follows $T$ through $D$ to the same socket surface. The rows $H$ and $C$ transport and replay that finite comparison, and $P,N$ record provenance and local naming.

For every stronger request, the packet has no coordinate to answer it. In particular, no row proves any of the four discharge hypotheses or exports a closed subject-reduction theorem. The consumer therefore remains confined to finite route agreement.
\end{proof}

\closureat{SubjectReductionRouteEquivalenceUp}{seedStr}

\begin{closurestatus}{\SubjectReductionRouteEquivalenceUp}
  \constructivestory{A $\SubjectReductionRouteEquivalenceUp$ packet is generated from bundle, consumption, componentwise $\hsame$, displayed $\Cont$, $\Pkg$, and local $\NameCert$ rows over the subject-reduction discharge socket.}
  \theoryclosure{\seedClosure}
  \scopeclosed{\autoref{def:subject-reduction-route-equivalence-carrier}, \autoref{def:subject-reduction-route-equivalence-classifier}, \autoref{thm:subject-reduction-route-equivalence-bundle-socket-handoff}, \autoref{thm:subject-reduction-route-equivalence-namecert-obligations}, and \autoref{thm:subject-reduction-route-equivalence-consumer-nonescape} bind finite route agreement for the bundle subject-reduction discharge route.}
  \formalstatus{\unformalizedV}
  \bridgestatus{none}
  \origin{human}
  \notclaimed{This chapter does not inhabit the discharge bundle, provide a closed subject-reduction theorem, verify the packet in Lean, or bridge it to a public schema.}
  \upgradepath{Obligation closure requires checked carrier admission, classifier exactness, socket handoff, NameCert obligations, consumer non-escape, provenance, and local naming for the same finite route-equivalence packet.}
\end{closurestatus}
